% zref-savepos based markers (three-pass for accurate page numbers)
% We record positions with zref-savepos and emit NDJSON lines for the Node builder.
% For accurate page numbers, especially for floats, we use zref's page reference
% system which requires a minimum of two passes, sometimes three for complex layouts.
% The coordinates from zposx/zposy are always from the previous run.
%
% IMPORTANT: Run LaTeX at least twice, preferably three times for documents
% with floats to ensure accurate page numbers in the output NDJSON.

% Ensure macros exist even if the package is missing (no-ops)
\providecommand\zsavepos[1]{}
\providecommand\zposx[1]{0}
\providecommand\zposy[1]{0}
\providecommand\zpageref[1]{1}

\makeatletter
\providecommand\zref@extractdefault[3]{#3}
\providecommand\zref@addprop[2]{}
\providecommand\zref@refused[1]{}

% Configure zref to save page numbers in addition to positions
% This is needed for accurate float page detection
\@ifpackageloaded{zref-savepos}{%
  % Add 'page' property to the savepos property list
  \zref@addprop{savepos}{page}%
}{}

% Note: Column index calculation in \geomemit:
% - Floats (type=table, type=figure): Always col=0 (atomic units, both start/end same value)
% - Non-floats (type=para, etc.): Calculated from X position (more reliable than \if@firstcolumn)
% - Multi-column layout is auto-detected by JavaScript: textwidth > columnwidth * 1.5
%
% Element types: para, table, figure, section, title, etc.
% These types are explicitly marked in LaTeX and used by JavaScript for processing
\makeatother

% Open a per-job NDJSON file for TeX positions
\newwrite\geomout
\immediate\openout\geomout=\jobname-texpos.ndjson
\AtEndDocument{\immediate\closeout\geomout}

% Helper to emit a JSON line to both log and NDJSON file for the Node builder
\makeatletter
\def\geomemit#1#2#3{%% #1=id, #2=role, #3=type (para, table, figure, section, etc.)
  \begingroup
    \edef\gid{#1}%%
    \edef\grole{#2}%%
    \edef\gtype{#3}%%
    \edef\gx{\zposx{gm:#1:#2}}%% scaled points
    \edef\gy{\zposy{gm:#1:#2}}%% scaled points
    % Use zref@extractdefault to get ACTUAL page number from zref label
    % This requires 2-3 LaTeX runs but gives accurate page numbers for floats
    % The "page" property is automatically saved by zref when using \zsavepos
    \edef\gpage{\zref@extractdefault{gm:#1:#2}{page}{1}}%% actual PDF page from zref, default to 1 if not found
    \def\gpagesource{zref}%%
    \edef\gpw{\the\paperwidth}%% e.g., 597.50787pt for A4 with 1in margins
    \edef\gph{\the\paperheight}%%
    \edef\gcw{\number\dimexpr\columnwidth\relax}%% column width (sp)
    \edef\gtw{\number\dimexpr\textwidth\relax}%% text width (sp)
    \edef\gsep{\number\dimexpr\columnsep\relax}%% column separation (sp)
    % Calculate column index
    % For floats (table, figure), always use col=0 since they're atomic units
    % Both start and end markers should have the same column value
    \def\gcol{0}%% Default to column 0
    % Check if type is table or figure
    \def\geom@isfloat{0}%%
    \def\geom@checktype##1table##2\@nil{%
      \ifx\relax##2\relax\else\def\geom@isfloat{1}\fi
    }%%
    \expandafter\geom@checktype\gtype table\@nil
    \def\geom@checktype##1figure##2\@nil{%
      \ifx\relax##2\relax\else\def\geom@isfloat{1}\fi
    }%%
    \expandafter\geom@checktype\gtype figure\@nil
    % For non-floats (paragraphs, sections, etc.): calculate from X position
    \if\geom@isfloat0
      % Calculate boundary: columnWidth + columnSep/2
      \@tempdima=\columnwidth\relax
      \advance\@tempdima by 0.5\columnsep\relax
      \@tempdimb=\gx sp\relax
      \if@twocolumn
        \ifdim\@tempdimb>\@tempdima
          \def\gcol{1}%% Right column
        \else
          \def\gcol{0}%% Left column
        \fi
      \fi
    \fi
    % Note: For floats, col is always 0 (atomic units, not split)
    % Note: twocolumn flag removed - multi-column is auto-detected from measurements (twsp > cwsp * 1.5)
    % Write NDJSON (no line wrapping) with source info for debugging
    \immediate\write\geomout{ {"id":"\gid","role":"\grole","type":"\gtype","xsp":"\gx","ysp":"\gy","pw":"\gpw","ph":"\gph","page":\gpage,"page_source":"\gpagesource","cwsp":\gcw,"twsp":\gtw,"col":\gcol,"colsep":\gsep} }%%
    % Also echo to log for debugging (may wrap)
    \message{GEOM: {"id":"\gid","role":"\grole","type":"\gtype","xsp":"\gx","ysp":"\gy","pw":"\gpw","ph":"\gph","page":\gpage,"page_source":"\gpagesource","cwsp":\gcw,"twsp":\gtw,"col":\gcol,"colsep":\gsep}}%% NDJSON duplicate
  \endgroup
}
\makeatother

% Inline mark attaches to the current line; safe in both paragraph and vertical mode
% #1=id, #2=role, #3=type
% DISABLED: Using line-level marking instead (see geom-marks-lines.tex)
\def\geommarkinline#1#2#3{%
  % No-op: Line-level marking provides better accuracy
}

% Immediate mark for floats (figures, tables) - improved version
% This version uses \label mechanism to get more accurate page numbers
% #1=id, #2=role, #3=type
\def\geommarknow#1#2#3{%
  \zsavepos{gm:#1:#2}%
  \label{geom:#1:#2}% Create a label for cross-referencing
  \geomemit{#1}{#2}{#3}%
}

% Advanced float mark that attempts to get the actual float page number
% This should be used within float environments for maximum accuracy
% #1=id, #2=role, #3=type
\def\geommarkfloat#1#2#3{%
  \zsavepos{gm:#1:#2}%
  \label{geom:#1:#2}%
  % For now, emit immediately to avoid macro expansion issues
  % TODO: Re-enable afterpage when label extraction is fixed
  \geomemit{#1}{#2}{#3}%
}

% ---------------------------------------------------------------------------
% LuaTeX Line-Level Marking Support
% ---------------------------------------------------------------------------
\ifx\directlua\undefined
  \message{GEOM: LuaTeX not detected. Line-level marking disabled.}
  \def\geommarkpara#1{}
\else
  % Define Lua functions first
  \directlua{
    geom_line_output_filename = tex.jobname .. "-texpos-lines.ndjson"
    geom_current_para_id = nil
    
    -- Clear/create fresh line-level NDJSON file at initialization
    local f = io.open(geom_line_output_filename, "w")
    if f then f:close() end
    
    function geom_write_line_data(id, line, w, h, d)
        local x = tex.lastxpos or 0
        local y = tex.lastypos or 0
        local page = tex.count[0] or 1
        
        local json = string.char(123) .. '"id":"' .. id .. '","role":"P","type":"line","line":' .. line .. ',"xsp":"' .. x .. '","ysp":"' .. y .. '","page":' .. page .. ',"w":' .. w .. ',"h":' .. h .. ',"d":' .. d .. string.char(125) .. string.char(10)
        
        local f = io.open(geom_line_output_filename, "a")
        if f then
            f:write(json)
            f:close()
        end
    end

    function geom_linebreak_callback(head)
        if not geom_current_para_id then return head end
        
        local para_id = geom_current_para_id
        geom_current_para_id = nil
        
        for n in node.traverse_id(node.id("hlist"), head) do
            local line_count = (n.geom_line_number or 0) + 1
            n.geom_line_number = line_count
            
            local save_pos = node.new("whatsit", "save_pos")
            local latelua = node.new("whatsit", "late_lua")
            latelua.data = 'geom_write_line_data("' .. para_id .. '", ' .. line_count .. ', ' .. n.width .. ', ' .. n.height .. ', ' .. n.depth .. ')'
            
            if n.list then
                n.list = node.insert_before(n.list, n.list, save_pos)
                node.insert_after(n.list, save_pos, latelua)
            end
        end
        
        return head
    end
    
    texio.write_nl("GEOM: Lua functions defined")
  }%
  
  % Line-level marking moved to geom-marks-lines.tex
  % This avoids conflicts with Elsevier template callbacks
  \def\geommarkpara#1{%
    % No-op: Load geom-marks-lines.tex for line-level marking
  }%
\fi